\section{Historical background}
\label{sec:history}

The geometry we use is ancient. Apollonius is credited with showing that the
motion of a body on an eccentric circle is equivalent to motion on an
epicycle carried by a concentric deferent \cite{Neugebauer1975}, and
Hipparchus is the traditional source of the quantitative use of these
devices for the Sun and the Moon \cite{Neugebauer1975,Evans1998}. In the
\emph{Almagest}, Ptolemy combined deferent, epicycle and eccentric for
every planet, and introduced the equant, a point offset from the deferent's centre,
opposite the Earth, about which the epicycle's centre, moving along the
deferent, sweeps out equal angles in equal times
\cite{Toomer1998,Neugebauer1975}. The equant is not a harmonic device;
Evans showed that it reproduces Kepler's second law to first order in the
eccentricity \cite{Evans1984}, and the later history of the division of the
Martian eccentricity can be read as a sequence of approximations to Kepler
motion \cite{Evans1988}.

Copernicus changed the origin of the frame but not the circle paradigm: his
models still need epicycles, with epicyclets replacing the equant
\cite{SwerdlowNeugebauer1984}. The step that makes the geocentric motion a
\emph{difference} of two heliocentric motions is already visible for
circular orbits, where the Earth's orbit becomes the epicycle of a superior
planet \cite{Evans1998,Scholz2022}. Kepler's ellipses then replaced the
circles \cite{Boccaletti2001}, yet an ellipse is not an escape from
circular motion: expanding Kepler's equation in a Fourier series in the
mean anomaly leads to Bessel-function coefficients, a result usually dated to
Lagrange and Bessel in the early nineteenth century \cite{Dutka1995}. In
this reading the ellipse is an infinite sum of uniform circular motions,
with coefficients that decay geometrically
(Sec.~\ref{sec:th-conv}). The connection between epicycles and
quasi-periodic expansions was drawn explicitly by Gallavotti
\cite{Gallavotti2001}, and Sec.~\ref{sec:prior} places our construction
relative to that literature. We cite these historical works through
secondary sources throughout and have not consulted the primary texts of
Bessel or Apollonius directly.
